\documentclass[11pt]{article}
\usepackage[margin=1in]{geometry}
\usepackage{amsmath,amssymb}
\usepackage{lmodern}
\usepackage{xurl}
\setlength{\parskip}{0.5em}
\setlength{\parindent}{0pt}
\begin{document}

\begin{center}
{\large\bfseries Anatomy of a Smooth Transition: A Task-Based Macroeconomic
Model of AGI Displacement with Fiscal and Social Safety Nets}

\vspace{0.8em}

DUAN Fang$^{*}$

\vspace{0.4em}

\emph{College of Information Science and Engineering, Huaqiao University,
Xiamen, China}

\vspace{0.2em}

{\small duan.fang@hqu.edu.cn}

\vspace{0.3em}

{\small $^{*}$Corresponding author.}
\end{center}

\vspace{0.6em}

\noindent\textbf{Abstract.}
As a transparent thought experiment, this paper models an economy where
artificial general intelligence (AGI) enters new-task creation itself, so the
engine reinstating displaced labor can collapse endogenously. We build a
two-layer task model whose reinstatement rate declines as AI capability
advances, with nominal rigidity, the zero lower bound and debt-deflation,
capital-tax financing, and a social-instability channel independent by
construction. Under full substitution, automation keeps exceeding labor's
capture of new tasks, driving market employment monotonically to zero without
self-correction; a high market-employment endpoint is inadmissible; it is the
counterfactual where full substitution is not reached. The piecewise-smooth
fast subsystem exhibits a persistence-type boundary equilibrium bifurcation
with a closed-form critical reinstatement rate, while the global outcome is a
different bifurcation. With public employment as full absorption under the
within-period fiscal constraint, the model admits two basins: a crash, or a
smooth transition with zero residual unemployment and consumption above the
floor, separated by two explicitly modeled gates, a fiscal gate on the
consumption floor and a social gate combining a social-instability threshold
with an adjustment-speed floor. A bisection locates the fiscal boundary at
$g_A=0.097$ under $\tau=0.15$; at higher tax rates consumption stays above the
floor at every admissible speed and the operative boundary is the
adjustment-speed gate near $g_A=0.14$, a minimum transition time rather than a
Laffer limit. A same-budget comparison that reserves revenue for genuine
positive transfer flows contracts the frontier, whereas the workfare-first
transfer buffer does not move it. On the smooth path at $g_A=0.010$, minimum
consumption is $0.415$, public employment ends at $0.85$ and residual
unemployment is zero, while the fixed corner holds $0.357$ and the no-policy
path crashes at $0.057$. Three magnitudes use U.S. data and the rest are
declared settings. A smooth transition therefore requires both a financed
fiscal base and a transition speed society can absorb, given automation speed.

\noindent\textbf{Keywords:} Artificial general intelligence; Task-based model;
Boundary equilibrium bifurcation; Geometric singular perturbation theory;
Public employment; Labor share.

\noindent\textbf{JEL classification:} E24, E31, E62, J64, O33.

\noindent\textbf{Acknowledgments.} The author thanks Mr. Guangtong Lu and
Mr. Xingbin Lai for valuable discussions and assistance.

\end{document}
